% Canonical extended proofs for ramified-gauss-core.tex.
% DR2109.07463v3 PDF df888f8527b3647a3428d6677a3e3070aab7439f79f737f3e894218fb64cfef0.
% de Faveri2610.04045v1 PDF 97720b3cacccbedb9c2014520ad4d128ae23afd0a781bdeff27dab4e86b7c726.
% Full primary bodies and private reading caches are excluded from public closure.
\section{Proofs of the ramified transform and full-coefficient prefix bound}
\label{rg:details}
The external theta input is the unconditional Section~5 and Appendix~A
of \href{https://arxiv.org/abs/2109.07463v3}{Dunn--Radziwi\l\l,
\emph{Bias in cubic Gauss sums: Patterson's conjecture},
arXiv:2109.07463v3, 14 May 2024}.
Its GRH-conditional prime theorems are not used.
The matrix input is the already admitted physical ideal family in
\href{https://arxiv.org/abs/2610.04045v1}{de Faveri,
\emph{Optimal large sieve for fixed order characters},
arXiv:2610.04045v1, Theorem~1.1 and Example~2.2}, at order three.
The theta modularity and deep cubic norm remain external theorems.
The constructions and estimates below specify how they apply to the
actual finite-character family.

\subsection{Fixed residue data, angular order and raw evaluations}
\label{rg:finite-family-proof}
The full varying-column phase list, after the arithmetic
\(\mu A_U\) opening, consists of one \(\vartheta\) or its conjugate,
the cubic Gauss root, the original finite-order \(\nu\), finite
source/ray factors, the mod-\(4\) Fourier factors of \(\Gamma\),
the actual primitive \(S\)-type factors, and residue symbols at the
fixed numerators \(u_0,\gamma_S,\mathfrak e,\mathfrak w\).
The set \(N\mathfrak w<64\), the \(S\)-types and the \(\mathfrak e\)
set are fixed finite sets. The growing \(\mathfrak k\) numerator
belongs to the cubic frequency instead of this list.

There is a common fixed residue modulus for this list.
A finite-order Hecke character has infinity type zero and factors
through its actual conductor. The source group is explicitly a ray
quotient: Section~2.3.2 of de Faveri defines
\(\mathcal C=H_{\mathfrak c}\otimes\mathbb Z/6\mathbb Z\)
for one sufficiently deep \(S\)-supported ray modulus \(\mathfrak c\).
The chosen primary generators realize these ray characters as finite
residue characters. The \(\Gamma\) factors are characters of
\((\mathcal O/(4))^\times\).

The local sixth-power tests also have a fixed finite modulus.
For a prime \(\mathfrak p\), Hensel's criterion at \(1\) gives
\(1+\mathfrak p^e\subset(K_{\mathfrak p}^\times)^6\)
when \(e>2v_{\mathfrak p}(6)\).
One may use \(e_{(2)}=3\) and \(e_{(\lambda)}=5\).
These are moduli for testing local classes, not replacements for
primitive character evaluation moduli.
For a fixed nonzero integral numerator \(\zeta\), its
\(n\)-th residue function, \(n\mid6\), is a finite Kummer character
for \(K(\zeta^{1/n})/K\).
A safe sextic conductor bound is supplied by the polynomial
discriminant ideal \((6)^6(\zeta)^5\).
A field/order discriminant divides this separable polynomial
discriminant, including when a proper factor supplies the field.
A fractional numerator can be cleared by a fixed sixth power,
retaining its original numerator and denominator exclusions.
This uses finite one-dimensional Kummer characters, without a
general Artin-holomorphy hypothesis.

A joint finite \(S\)-type or primitive-value selector is not assumed
multiplicative. If \(\mathcal H\) is its fixed finite residue group,
write its actual function as
\begin{equation}\label{rg:joint-selector}
 \begin{aligned}
 f(x,y)&=\sum_{\chi,\psi\in\widehat{\mathcal H}}
                           b_{\chi,\psi}\chi(x)\psi(y),\\
 b_{\chi,\psi}
   &=|\mathcal H|^{-2}\sum_{x,y\in\mathcal H}
                   f(x,y)\overline{\chi(x)\psi(y)},\\
 \sum_{\chi,\psi}|b_{\chi,\psi}|
   &\le|\mathcal H|\,\|f\|_\infty .
 \end{aligned}
\end{equation}
Normalized Parseval and Cauchy prove the last line.
This expansion precedes freezing a primitive type or separating
column factors. Forbidden primitive \(\mathfrak e/\mathfrak w\)
values are zero in \(f\); individual Fourier terms can be nonzero
there, and only their full sum reconstructs the zero.
No term is individually renamed an original nonzero row.

Thus one may use a common fixed raw modulus \(F_{\rm raw}\) and
functions
\begin{equation}\label{rg:raw-character}
 \kappa_{\rm raw}(\mathfrak a)
 =\kappa_{\rm prim}(\alpha_{\mathfrak a})
                         \mathbf1_{(\mathfrak a,F_{\rm raw})=1}.
\end{equation}
They are imprimitive residue characters on that modulus.
All stipulated prime zeros survive a primitive conductor drop.

The Jacobi function \(\rho_J=\mu\vartheta\) is not one of these
finite characters. Its \(\mu\) has already entered the arithmetic
coefficient \(B_U\), and its unique remaining angular factor is
\(\vartheta\), with order \(+1\) or \(-1\) in the two conjugate
orientations. Common or prefix Jacobi factors stay in their actual
outer coefficients when those variables are frozen.
An arbitrary bounded multiplicative function, a growing opposite
column, a M\"obius coefficient or the nonperiodic density
\(\prod_{\mathfrak p\mid\mathfrak a}(1-(N\mathfrak p)^{-1})\)
is not made into a member of this fixed family. That density does
not occur in \(\mathcal P_X\).

\subsection{Known integral cusps and an explicit absolute coefficient norm}
\label{rg:cusp-norm-proof}
Let
\(\Gamma=\mathrm{SL}_2(\mathcal O)\),
\(\Gamma_1(3)=\ker(\Gamma\to\mathrm{SL}_2(\mathcal O/(3)))\)
and
\(\Gamma_2=\mathrm{SL}_2(\mathbb Z)\Gamma_1(3)\).
The source theta function is a weight-zero function on hyperbolic
three-space with
\(\theta(\delta w)=\chi(\delta)\theta(w)\) for
\(\delta\in\Gamma_2\).
The multiplier \(\chi\) is the source Kubota character.
The quotient \(\Gamma_2\backslash\Gamma\) has \(27\) cosets.
One can see the finite count from
\(\mathcal O/(3)\simeq\mathbb F_3[\varepsilon]/(\varepsilon^2)\):
the reduction kernel in \(\mathrm{SL}_2\) has
\(3^3\) elements, and \(\mathrm{SL}_2(\mathbb F_3)\) has \(24\).

Fix the published integral representatives \(\gamma_j\) and put
\(F_j(w)=\theta(\gamma_jw)\).
For the exact chosen function define its actual Fourier coefficient
\(a_j(\mu)\) by
\begin{equation}\label{rg:cusp-expansion}
 F_j(z,v)=\sigma_jv^{2/3}
  +\sum_{\mu\ne0}a_j(\mu)vK_{1/3}(4\pi|\mu|v)
                                        e(\operatorname{Tr}(\mu z)),
 \qquad e(x)=e^{2\pi ix}.
\end{equation}
For those representatives
\(\sigma_j=3^{5/2}/2\) for \(j\le9\) and zero for \(j\ge10\).
The support lies in \(\lambda^{-4}\mathcal O\).
The symbol \(a_j\) always means this actual coefficient;
a source table using \(d_j(\mu)\), \(d_j(-\mu)\) or a conjugated
function is converted before using an exact phase formula.

The source formulas (5.7), (5.13), (5.14) and Appendix~A express
these coefficients through \(\tau,\tau_1,\tau_2\), unit rotations,
fixed unit factors and additive phases.
The first nine families have the \(\tau\) absolute norm and the
other eighteen have the \(\tau_1,\tau_2\) norms.
Rotations, conjugation and index reversal preserve absolute norm.

\begin{proposition}[Known-cusp absolute norm]\label{rg:cusp-norm}
For \(1<\sigma<13/12\), every one of these actual cusp families satisfies
\begin{equation}\label{rg:cusp-norm-bound}
 \mathcal A_j(\sigma):=
   \sum_{\mu\ne0}\frac{|a_j(\mu)|}{(N\mu)^\sigma}
 \le3^8\zeta_K^{(\lambda)}(\sigma)
                      \zeta_K^{(\lambda)}(3\sigma-\tfrac12).
\end{equation}
The right side is independent of a moving good level or a translation.
Here \(\zeta_K^{(\lambda)}\) omits the prime above \(3\).
\end{proposition}
\begin{proof}
In the \(\tau\) formula write \(c,d\equiv1\pmod3\), with \(c\)
squarefree and \(d\) arbitrary, possibly sharing primes with \(c\).
For the first family
\(\mu=\pm\epsilon\lambda^{3n-4}cd^3\), \(n\ge1\),
there are at most six sign/unit branches and
\[
 |\tau(\mu)|\le3^{n/2+2}\sqrt{Nd}.
\]
For the second family
\(\mu=\pm\lambda^{3n-3}cd^3\), \(n\ge0\),
there are at most two branches and
\(|\tau(\mu)|\le3^{n/2+5/2}\sqrt{Nd}\).
These follow from the literal unnormalized primitive Gauss magnitude
\(|g(a,c)|=\sqrt{Nc}\) and the source factor \(|d/c|\);
the fixed Gauss numerators are units modulo good \(c\).
Retain the separate \(c\) and cube variables. A generic
\((N\mu)^{1/6}\) coefficient bound would lose convergence near \(1\).

Put \(z=3^{1/2-3\sigma}\) and
\(Z_{\rm SF}(\sigma)=\sum_{c\ {\rm primary,SF}}(Nc)^{-\sigma}\),
with the prime \(\lambda\) excluded. Then
\begin{equation}\label{rg:tau-geometric}
 \begin{split}
 \mathcal A_\tau(\sigma)\le
 \left\{\frac{6\,3^{2+4\sigma}z+
                         2\,3^{5/2+3\sigma}}{1-z}\right\}
 Z_{\rm SF}(\sigma)\zeta_K^{(\lambda)}(3\sigma-\tfrac12).
 \end{split}
\end{equation}
The \(\tau_1,\tau_2\) families each have three unit branches at
\(\lambda^{-4}cd^3\), with magnitude at most \(9\sqrt{Nd}\);
their possible opposite signs do not change the norm. Consequently
\begin{equation}\label{rg:tau12-geometric}
 \mathcal A_{\tau_i}(\sigma)
 \le27\,3^{4\sigma} Z_{\rm SF}(\sigma)
                    \zeta_K^{(\lambda)}(3\sigma-\tfrac12),
 \qquad i=1,2 .
\end{equation}
In the stated range \(z\le1/9\),
\(3^{2+4\sigma}\le3^7\),
\(3^{5/2+3\sigma}\le3^6\), and
\(27\,3^{4\sigma}\le3^8\).
The brace in \eqref{rg:tau-geometric} is at most \(3^8/2\).
Finally \(Z_{\rm SF}(\sigma)\le\zeta_K^{(\lambda)}(\sigma)\).
The actual coefficient convention and all unit phases therefore
give \eqref{rg:cusp-norm-bound}.
\end{proof}

For later comparison, an affine state
\(F_j(Az+B,|A|v)\), \(A\ne0\), has index and coefficient
\begin{equation}\label{rg:affine-coefficient}
 \nu=A\mu,\qquad
 a_{j,A,B}(\nu)=|A|\,a_j(\mu)e(\operatorname{Tr}(\mu B)).
\end{equation}
Its absolute Dirichlet norm is exactly
\(|A|^{1-2\sigma}\mathcal A_j(\sigma)\).
For the ramified diagonal \(A=\lambda^{-2B_0}\), this factor is
\(3^{B_0(2\sigma-1)}\).
Commuting a moving translation past that diagonal changes
\(U(-\lambda^4c/r)\) to \(U(-\lambda^{4+2B_0}c/r)\).
Together with \(\nu=\lambda^{-2B_0}\mu\) its phase is still
\(e(-\operatorname{Tr}(\lambda^4\mu c/r))\).
This explains why an omitted ramified scale changes actual phases,
even though the eventual absolute bound is fixed-data controlled.

\subsection{Exact rational reduction and the actual Kubota multiplier}
\label{rg:rational-construction}
For \(a=A/C\in K\) reduced, choose \(B_0,D_0\in\mathcal O\) with
\(AD_0-B_0C=1\) and put
\(\gamma_a=\left(\begin{smallmatrix}A&B_0\\C&D_0\end{smallmatrix}\right)\).
The Euclidean property of \(\mathcal O\) supplies this choice.
Let
\[
 U(x)=\begin{pmatrix}1&x\\0&1\end{pmatrix},\qquad
 W_C=\begin{pmatrix}0&-1/C\\C&0\end{pmatrix}.
\]
Integral \(a\) is included by \(C=1\).
Direct multiplication proves
\begin{equation}\label{rg:rational-matrix}
 \gamma_a^{-1}U(a)
 =\begin{pmatrix}D_0&1/C\\-C&0\end{pmatrix}
 =U(-D_0/C)W_{-C}.
\end{equation}
Choose \(j\) with \(\delta_a=\gamma_a\gamma_j^{-1}\in\Gamma_2\).
Then the exact function identity is
\begin{equation}\label{rg:rational-theta}
 \theta(U(a)w)
 =\chi(\delta_a)F_j(U(-D_0/C)W_{-C}w).
\end{equation}
It uses known integral cusp functions rather than an unconstructed
new cusp coefficient family.

The phase is computable, not an unspecified unit set to one.
Find a rational integral \(\sigma\in\mathrm{SL}_2(\mathbb Z)\)
with the same reduction modulo \(3\) as \(\delta_a\).
Then \(\delta'=\sigma^{-1}\delta_a\in\Gamma_1(3)\).
If \(\delta'=\left(\begin{smallmatrix}a'&b'\\c'&d'\end{smallmatrix}\right)\),
the source rule is
\begin{equation}\label{rg:kubota-rule}
 \chi(\delta_a)=\chi(\delta')=
 \begin{cases}
  (c'/a')_3,&c'\ne0,\\
  1,&c'=0 .
 \end{cases}
\end{equation}
Here \(a'\equiv1\pmod3\) is nonzero and coprime to \(c'\).
The index \(j\) can be found by reducing modulo \(3\) and testing the
\(24\)-element rational subgroup against the \(27\) fixed cosets.
Equivalently, a finite traversal of elementary generators constructs
representatives. A representative differing from the published one
changes its function by the fixed multiplier just computed.
Its Fourier coefficients and constants carry that multiplier too.

One can also display the fixed ramified states explicitly.
If a moving good level \(r\) survives the fixed denominator reduction,
write \(C=c_0r\), with \(c_0\) in a fixed finite set and
\((c_0,r)=1\). The matrix identity
\[
 U(A/C)W_r=\gamma_a U(-D_0/C)
                       \operatorname{diag}(r/C,C/r)
\]
and \(D_0=br+9c_0d\) give the fixed state
\(F_j|U(-b/c_0)\operatorname{diag}(c_0^{-1},c_0)\)
followed by the moving translation \(U(-9c_0^2d/r)\).
The general formula \eqref{rg:rational-theta} also handles cancelled
or nonunit fixed Fourier residues directly; this special display
does not impose a primitive Fourier-support condition on them.

\subsection{Radial cusp constants, angular derivatives and the Mellin Jacobian}
\label{rg:mellin-proof}
Write \(w=(z,v)\). The action in \eqref{rg:rational-theta} is
\begin{equation}\label{rg:cusp-action}
 z'=-D_0/C-\frac{\bar z}{C^2(|z|^2+v^2)},\qquad
 v'=\frac{v}{NC(|z|^2+v^2)} .
\end{equation}
The pulled-back constant of \(F_j\) is therefore
\begin{equation}\label{rg:pulled-constant}
 \chi(\delta_a)\sigma_j(NC)^{-2/3}
                  v^{2/3}(|z|^2+v^2)^{-2/3}.
\end{equation}
It is radial in \(z\). Every nonzero pure Wirtinger derivative
\(\partial_z^m\) or \(\partial_{\bar z}^m\) at \(z=0\) kills it.
At large \(v\) the original translated theta constant is also
independent of \(z\). The known nonconstant lattice support gives
exponential decay after angular differentiation at both ends,
with the reciprocal-end rate depending on \(NC\).

For completeness, the zeroth-angular constants have not disappeared.
For a finite combination \(\sum_a c(a)\theta(z+a,v)\) put
\begin{equation}\label{rg:endpoint-constants}
 C_\infty=\frac{3^{5/2}}2\sum_a c(a),\qquad
 C_0^*=\sum_a c(a)\chi(\delta_a)\sigma_{j(a)}(NC_a)^{-2/3}.
\end{equation}
Its regularized Mellin transform with weight \(v^{2s-2}\) has residues
\begin{equation}\label{rg:zeroth-residues}
 \operatorname*{Res}_{s=1/6}\Lambda(s,0)=-C_\infty/2,
 \qquad
 \operatorname*{Res}_{s=5/6}\Lambda(s,0)=C_0^*/2 .
\end{equation}
These signs follow by integrating the powers at infinity and zero.
A projector with zero sum kills \(C_\infty\), but need not kill
\(C_0^*\). No finite noncubic label implies such a cancellation.

For \(\ell>0\) set
\(\mathcal D_\ell=(2\pi i)^{-\ell}\partial_z^\ell\);
for \(\ell<0\) use
\(\mathcal D_\ell=(2\pi i)^{-|\ell|}\partial_{\bar z}^{|\ell|}\).
The angular Mellin transform is
\begin{equation}\label{rg:angular-mellin}
 \Lambda_a(s,\ell)=\int_0^\infty
  \mathcal D_\ell[\theta(z+a,v)]_{z=0}
                                v^{2s+|\ell|-2}\,dv .
\end{equation}
For \(\ell\ne0\) it is entire by the endpoint bounds just proved.

\begin{proposition}[Exact angular rational-cusp identity]
\label{rg:angular-identity}
For every reduced rational translate and \(\ell\ne0\),
\begin{equation}\label{rg:mellin-identity}
 \Lambda_a(s,\ell)=
 \chi(\delta_a)(-1)^{|\ell|}(C/\bar C)^\ell(NC)^{1-2s}
                   \Lambda_{j,-D_0/C}(1-s,-\ell).
\end{equation}
The right Mellin transform uses the actual known function
\(F_j(z-D_0/C,v)\).
\end{proposition}
\begin{proof}
For \(\ell>0\), set \(\bar z=0\) in the full-coordinate chain rule.
Then \(z'\) and \(v'\) in \eqref{rg:cusp-action} are constant as
functions of \(z\), while \(\bar z'\) is affine with derivative
\(-1/(\bar C^2v^2)\).
Thus the pure derivative has exactly the factor
\(\chi(\delta_a)(-1)^\ell\bar C^{-2\ell}v^{-2\ell}\)
and the opposite angular derivative of \(F_j\) at
\((-D_0/C,1/(NCv))\).
There are no lower angular or radial terms.
The negative sign case is the conjugate-coordinate calculation.

In the integral substitute \(y=1/(NCv)\).
Its Jacobian and power weight give
\((NC)^{-2s+\ell+1}y^{\ell-2s}\,dy\).
Combining the factor \(\bar C^{-2\ell}\) produces precisely
\((C/\bar C)^\ell(NC)^{1-2s}\).
The exponent \(1-2s\), rather than \(-2s\), is essential.
This proves \eqref{rg:mellin-identity} for both signs.
\end{proof}

If the real part of the dual Dirichlet variable exceeds one, its series is
\begin{equation}\label{rg:dual-translated-series}
 \sum_{\mu\ne0}a_j(\mu)e(-\operatorname{Tr}(\mu D_0/C))
       (\mu/|\mu|)^{-\ell}(N\mu)^{-(1-s)} .
\end{equation}
The additive and angular phases remain literal; their absolute values
are one only after this exact expression has been established.
Proposition~\ref{rg:cusp-norm} controls the series.

\subsection{The integral-index primary projector and raw cube completion}
\label{rg:projector-proof}
Let \(\kappa\) be any one fixed residue member on the primary section.
Use a fixed modulus containing its actual conductor and all original
fixed exclusions. Separate its \(\lambda\)-part \(\kappa_\lambda\)
and away-\(\lambda\) part \(\kappa_0\).
Choose \(f\ge2\) containing the \(\lambda\)-conductor, and put
\begin{equation}\label{rg:primary-projector}
 \Phi(x)=\kappa_\lambda(x)
               \mathbf1_{x\equiv1\bmod\lambda^2}
          \quad(x\bmod\lambda^f),
 \qquad \Phi(x)=0\text{ on nonunits}.
\end{equation}
The entire theta support has \(\lambda^3\mu\in\mathcal O\).
The correct argument is this integral index, not
\(\lambda^4\mu\) for a ramified character.

From the two \(\tau\) families above, \(\Phi(\lambda^3\mu)\)
selects exactly
\begin{equation}\label{rg:primary-support}
 \mathcal S=\{\lambda^{-3}cd^3:
               c,d\equiv1\pmod3,\ c\ {\rm squarefree}\}.
\end{equation}
Indeed the first family has positive \(\lambda\)-valuation after
multiplication by \(\lambda^3\), as do the second-family terms
with \(n\ge1\). They are killed as nonunits.
At \(n=0\) the negative branch is \(-1\) modulo \(3\),
and only the positive branch is \(1\).
This proves the unit-state selection as well as its valuation condition.

For a moving good squarefree level \(\mathfrak t\), let
\(\psi_{\mathfrak t}\) be either cubic orientation on its literal
residue modulus. Let \(F_0\) be the fixed away-\(\lambda\) raw modulus.
The product
\(\psi_M(x)=\kappa_0(x)\psi_{\mathfrak t}(x)\)
is a bounded periodic function on the valid period
\(M=F_0\mathfrak t\).
This period is not asserted to be its primitive conductor.
It is allowed to be imprimitive and to have overlapping
\(F_0,\mathfrak t\). No value \(\kappa(\mathfrak t)\) is divided out.

Define the coefficient-projected theta function using
\(\Phi(\lambda^3\mu)\psi_M(\lambda^3\mu)\).
Its infinite constant is zero. On \(\mathcal S\),
\(\tau(\mu)=3^{5/2}\widetilde g(c)\sqrt{Nd}\).
Multiplicativity on the primary section gives the exact cube identity
\begin{equation}\label{rg:cube-coefficients}
 \Phi(\lambda^3\mu)\psi_M(\lambda^3\mu)
 =\kappa(c)\kappa(d)^3\psi_{\mathfrak t}(c)
                                    \mathbf1_{(d,\mathfrak t)=1}.
\end{equation}
Write
\begin{equation}\label{rg:gauss-dirichlet}
 \mathcal R_{\kappa,\mathfrak t,\ell}(s)
 =\sum_{c\ {\rm primary,SF}}
 \frac{\widetilde g(c)\kappa(c)\psi_{\mathfrak t}(c)
                                  \Theta_\ell(c)}{(Nc)^s},
\end{equation}
with every fixed raw zero in \(\kappa\). Then initially \(\Re s>1\),
\begin{equation}\label{rg:cube-completion}
 D_{\Phi,\psi}(s,\ell)
 =3^{5/2+3s}i^\ell
 L^{\rm raw}(3s-\tfrac12,\kappa^3\Theta_{3\ell})
                       \mathcal R_{\kappa,\mathfrak t,\ell}(s).
\end{equation}
The factor \(i^\ell\) is the actual angle of \(\lambda^{-3}\).
The cube series has denominator \((Nd)^{3s-1/2}\), and
\(L^{\rm raw}\) retains all original exclusions above
\(\lambda F_0\mathfrak t\), even after primitive conductor drop.
There is no \(\kappa(\lambda)\) factor or untwisted-zeta substitution.
Conjugating the whole construction gives the actual conjugate
Gauss root and opposite angular/cubic orientations.

The full ramified projector has an exact finite translation formula.
The trace dual of \(\mathcal O\) is \(\lambda^{-1}\mathcal O\), so
\begin{equation}\label{rg:ramified-fourier}
 \begin{aligned}
 b_h&=3^{-f}\sum_{x\bmod\lambda^f}
         \Phi(x)e(-\operatorname{Tr}(xh/\lambda^{f+1})),\\
 \Phi(x)&=\sum_{h\bmod\lambda^f}
                  b_he(\operatorname{Tr}(xh/\lambda^{f+1})),&
 \sum_h|b_h|&\le3^{f/2}.
 \end{aligned}
\end{equation}
For the away period use normalized Fourier coefficients \(q_u\)
with positive inversion sign and pairing
\(e(\operatorname{Tr}(xu/\alpha_M))\).
Since \(M\) avoids \(\lambda\), that pairing is perfect, and
\begin{equation}\label{rg:away-fourier}
 \psi_M(x)=\sum_{u\bmod M}q_u e(\operatorname{Tr}(xu/\alpha_M)),
 \qquad
 \sum_u|q_u|\le\sqrt{NM}.
\end{equation}
Both mass bounds are Parseval followed by Cauchy.
They hold for arbitrary bounded periodic functions; no primitive
Fourier-support condition is inserted.

Consequently the exact projected function is a finite sum of
\(\theta(z+a_{h,u},v)\) with
\begin{equation}\label{rg:actual-translates}
 a_{h,u}=h\lambda^{2-f}+\lambda^3u/\alpha_M,\qquad
 \sum_{h,u}|b_hq_u|\le3^{f/2}\sqrt{NM}.
\end{equation}
The inverse different explains the exponent \(2-f\).
Every raw zero is reconstructed by this full sum.
Each reduced denominator \(C_{h,u}\) divides a fixed
\(\lambda\)-power times \(M\); hence
\(1\le NC_{h,u}\ll_{\rm fixed}N\mathfrak t\).
Nonunit Fourier residues may lower that denominator. They remain
terms of the exact transform, rather than being discarded.

\subsection{The coarse dual bound, safe convexity and literal Perron endpoints}
\label{rg:sharp-proof}
Let \(R=N\mathfrak t\), and set
\[
 E(s)=L^{\rm raw}(3s-\tfrac12,\kappa^3\Theta_{3\ell})
                                  \mathcal R_{\kappa,\mathfrak t,\ell}(s).
\]
For \(\ell\ne0\) the finite sum of entire angular Mellins is entire.
The source Fourier--Mellin calculation has factor
\begin{equation}\label{rg:gamma-factor}
 \gamma_\ell(s)=\frac14(2\pi)^{-2s-|\ell|}
 \Gamma(s+\tfrac{|\ell|}{2}-\tfrac16)
 \Gamma(s+\tfrac{|\ell|}{2}+\tfrac16).
\end{equation}
Thus division by \(\gamma_\ell\), whose reciprocal is entire,
and the nonzero elementary factor in \eqref{rg:cube-completion}
makes \(E\) entire as well.

On \(s=-\eta+i\tau\), \(0<\eta<1/12\),
\eqref{rg:mellin-identity}, \eqref{rg:cusp-norm-bound} and
\eqref{rg:actual-translates} give
\begin{equation}\label{rg:dual-mass}
 \begin{split}
 \sum_{h,u}|b_hq_u|(NC_{h,u})^{1+2\eta}
                  \mathcal A_{j(h,u)}(1+\eta)
 &\ll_{\rm fixed,\eta} R^{3/2+2\eta},\\
 |E(-\eta+i\tau)|
 &\ll_{\rm fixed,\ell,\eta}
 R^{3/2+O(\eta)}(1+|\tau|)^{2+O(\eta)} .
 \end{split}
\end{equation}
The height exponent follows from the quotient of the two Gamma
products. The exact multiplier
\(\chi(\delta_{h,u})(-1)^{|\ell|}
 (C_{h,u}/\bar C_{h,u})^\ell\)
has been written before bounding its modulus.
This is a conservative \(R^{3/2+\epsilon}\) cost.
No refined \(R^{1/2+\epsilon}\) Ramanujan estimate is used.

The line \(\Re s=1+\eta\) is absolutely bounded uniformly in \(R\),
by the original Gauss and cube coefficients. For a requested final
\(\epsilon>0\), choose \(\eta>0\) sufficiently small in terms of
\(\epsilon\); the constants may depend on that choice.
The finite-order continuation from the theta Mellin representation
permits the damped strip maximum principle.
Applying it to the entire numerator, then allocating small exponents,
gives
\begin{equation}\label{rg:completed-convexity}
 |E(\sigma+i\tau)|
 \ll_{\rm fixed,\ell,\epsilon}
 (R(2+|\tau|))^\epsilon
 R^{\frac32(1-\sigma)}(1+|\tau|)^{2(1-\sigma)}
\end{equation}
on the safe half-strip \(1/2+\eta\le\sigma\le1+\eta\).
Only there divide by \(L^{\rm raw}\):
the real part of \(3s-\tfrac12\) exceeds one, so its inverse is
bounded by absolutely convergent Euler products, uniformly in all
raw omissions. No inverse-\(L\) continuation through unknown zeros
or zero-free strip is asserted.
Equation~\eqref{rg:completed-convexity} therefore bounds
\(\mathcal R\) on that half-strip.

Here are the endpoint details needed for
Theorem~\ref{rg:sharp-row}. Replace each real endpoint by the
half-integer threshold selecting exactly the same integer norms.
All ideals in a repeated norm fibre are counted. Such a fibre has
divisor-bounded multiplicity in this quadratic field.
For an unnormalized sharp prefix at length \(T\), truncated Perron
at height \(\mathcal H\ge1\) has error
\begin{equation}\label{rg:perron-error}
 O_\epsilon\bigl(T^\epsilon[1+(T/\mathcal H)\log(2T)]\bigr).
\end{equation}
This follows by summing
\(\min(1,T/(\mathcal H|n-y|))\) over integer distances from the
half-integer endpoint \(y\). No norm tie receives a half-weight.

At \(\sigma=1/2+\eta\) the vertical contribution is bounded by
\[
 (RT)^\epsilon T^{1/2}R^{3/4}\mathcal H .
\]
If \(T\ge R^{3/2}\), choose
\(\mathcal H=(T/R^{3/2})^{1/4}\).
On horizontal sides the bound is at most \(T/\mathcal H\),
up to allocated subpowers, since \(R^{3/2}\mathcal H^2\le T\).
The vertical integral, truncation and horizontal contributions
are then \(O((RT)^\epsilon T^{3/4}R^{3/8})\).
An interval is a difference of two prefixes. If \(T<R^{3/2}\),
the elementary \(O(T)\) count is the better branch.

Finally apply partial summation to the same
\(T^{-1/2}W(x)x^{-1/2}\), \(x=N\mathfrak a/T\).
Its supremum and total variation are bounded by
\(O(T^{-1/2}V_1(W))\). This proves
\eqref{rg:sharp-bound} with the exact profile.
For \(1/2\le T<1\) only the unit can occur; its coefficient and
profile are literal and satisfy the same bound with a fixed
endpoint constant. For \(T\le1/2\) there is no contribution.
This extension uses no analytic theorem at a fictitious length.

\subsection{The actual annular coefficient and exact eleven-factor expansion}
\label{rg:hb-proof}
The opened coefficient \(B_U=\mu A_U\) in
\eqref{rg:opened} is asserted on the actual annulus \(N\mathfrak m>U\).
It is distinct from original \(C_U\).
No identity for \(B_U(\mathcal O)\) is needed or asserted.
Unit factors within the factorization of an annular column are
essential and stay in every convolution below.

Let \(\mathbf1\) be the constant-one function on good ideals,
\(\delta_{\mathcal O}\) the ideal-convolution identity, and
\(M_z(\mathfrak a)=\mu(\mathfrak a)
 \mathbf1_{N\mathfrak a\le z_0}\), \(z_0=(2D)^{1/5}\).
The exact good-ideal convolution identity is
\begin{equation}\label{rg:hb-identity}
 \mu=\sum_{j=1}^5(-1)^{j-1}\binom5j
                   \mathbf1^{*(j-1)}*M_z^{*j}
       +\mathbf1^{*4}*(\mu-M_z)^{*5}.
\end{equation}
Indeed \(\mathbf1*\mu=\delta_{\mathcal O}\), so expand
\(\mu*[\delta_{\mathcal O}
       -(\delta_{\mathcal O}-\mathbf1*M_z)^{*5}]\);
also \(\delta_{\mathcal O}-\mathbf1*M_z
       =\mathbf1*(\mu-M_z)\).
Each factor of \(\mu-M_z\) has norm strictly above \(z_0\).
The last term in \eqref{rg:hb-identity} is supported strictly above
\(z_0^5=2D\) and vanishes at every residual
\(\mathfrak c\) under consideration.

The \(j\)-th term has \(j\) truncated M\"obius factors and \(j-1\)
unrestricted coefficient-one factors. Together with the two original
short factors it has \(2j+1\le11\) labels.
The \(j=1\) term and unit labels remain. Restricting the total column
to squarefree support forces all factors to be squarefree and
pairwise coprime.

Use independent dyads \([M_i,2M_i)\), \(M_i\ge1\), with the original
own-factor cutoffs at \(U,z_0\). If the box meets the annulus,
\begin{equation}\label{rg:box-product}
 D/2^{11}\le P_{\rm box}:=\prod_iM_i\le2D .
\end{equation}
Choose the partition from these scales only.
If some \(M_i\in[P_{\rm box}^{1/3},P_{\rm box}^{2/3}]\), it is
a balanced subset. If all \(M_i<P_{\rm box}^{1/3}\), multiply
successive scales until first reaching \(P_{\rm box}^{1/3}\);
the result is below \(P_{\rm box}^{2/3}\).
If no balanced subset exists, a unique scale exceeds
\(P_{\rm box}^{2/3}\); two such scales would have product above
\(P_{\rm box}\). The product of all remaining scales is below
\(P_{\rm box}^{1/3}\).

At \(D\ge2^{64}\),
\[
 P_{\rm box}^{2/3}\ge D^{2/3}/2^{22/3}>\sqrt D
                 \ge U,\qquad
 P_{\rm box}^{2/3}>(2D)^{1/5}.
\]
Therefore an unbalanced unique long factor is genuinely one of the
unrestricted coefficient-one factors. It cannot be an original
short factor or a truncated M\"obius factor.
The actual whole prefix \(\mathfrak p\), the product of all remaining
factors, satisfies \(N\mathfrak p\le2^{10}P_{\rm box}^{1/3}
 \ll D^{1/3}\).
The \(j=1\) term has no unrestricted factor and is automatically
balanced; its unit assignments have not been discarded.

For a fixed box define each macro coefficient by summing all its
signed assignments within its assigned factor group.
On a squarefree ideal the number of at most eleven labelled
factorizations is at most \(11^{\omega(\mathfrak a)}\).
Fixed binomial factors and dyadic logarithms thus give a uniform
\(D^\epsilon\) pointwise bound, with every own cutoff retained.
There is no separate count over \(N\mathfrak a,N\mathfrak b\le U\).

\subsection{Whole-factor cubic matrices and the large-prefix part}
\label{rg:matrix-proof}
A balanced partition gives two physical squarefree macro ideals
\(\mathfrak u,\mathfrak v\), each with scale in fixed multiples of
\([D^{1/3},D^{2/3}]\) and product comparable to \(D\).
Their own coefficients contain the original cutoffs and signs.
Cubic Gauss CRT gives on coprime products
\begin{equation}\label{rg:gauss-cross}
 \overline{G_{\mathfrak u\mathfrak v}^{\rm act}}
 =\overline{G_{\mathfrak u}^{\rm act}}\,
  \overline{G_{\mathfrak v}^{\rm act}}\,
                    \chi_{3,\mathfrak u}(\alpha_{\mathfrak v}).
\end{equation}
The raw matrix entry is zero at a collision, so summing all
squarefree macro pairs implements the original squarefree product.
No primitive root at an overlapping product is invented.

The external physical cubic norm is
\begin{equation}\label{rg:cubic-input}
 \bigl\|[\chi_{3,\mathfrak u}(\alpha_{\mathfrak v})]
                            _{N\mathfrak u\asymp P,\,
                              N\mathfrak v\asymp L}\bigr\|_{\rm op}
 \ll_\epsilon(PL)^\epsilon
                  \sqrt{P+L+(PL)^{2/3}} .
\end{equation}
The rows and columns are actual good squarefree ideals.
All raw fixed/growing multiplicative masks are coefficient
diagonals in this norm, after finite-selector expansion.
Ideal counting gives squared coefficient norms \(D^\epsilon\)
when each coefficient is weighted by the inverse square root norm.
Separate the one occurrence of the common \(W\) by
\[
 W(N\mathfrak uN\mathfrak v/D)
 =\frac1{2\pi}\int_{\mathbb R}\widehat W(t)D^{-it}
                       (N\mathfrak u)^{it}(N\mathfrak v)^{it}\,dt .
\]
The powers are unitary diagonals, costing \(\mathcal M(W)\) once.
Balanced boxes therefore cost
\(\mathcal M(W)D^{1/3+\epsilon}\).

For a Type-I box write \(\mathfrak m=\mathfrak p\mathfrak v\),
where \(\mathfrak v\) has coefficient one and the intrinsic
\(H_{\mathfrak p}\) contains all remaining assignments.
Then \(N\mathfrak p\ll D^{1/3}\).
The portion \(N\mathfrak p>Z\), with its literal own-variable
indicator, has dyadic bound
\[
 \mathcal M(W)D^\epsilon
       \sqrt{P+D/P+D^{2/3}},
\]
and hence total bound
\begin{equation}\label{rg:large-prefix}
 \mathcal M(W)D^\epsilon
                           [D^{1/3}+\sqrt{D/Z}].
\end{equation}
Macro own-factor constraints remain coefficients.
A chosen equivalent row-dependent norm interval costs the reviewed
logarithmic interval-Schur factor, preserving complete norm fibres.
It is not assigned operator norm one.

\subsection{The full original opposite-column mask in the small-prefix row}
\label{rg:deletion-proof}
Freeze the original \(\mathfrak n\), the exact
\(\mathfrak k,\mathfrak s,\mathfrak w,\mathfrak e\), and a finite
source component before estimating a Type-I row.
The original pair mask is exactly
\begin{equation}\label{rg:prefix-mask}
 \mathbf1_{(\mathfrak p\mathfrak v,\mathfrak n)=1}
 =\mathbf1_{(\mathfrak p,\mathfrak n)=1}
                 \mathbf1_{(\mathfrak v,\mathfrak n)=1}.
\end{equation}
The first factor stays in the actual prefix coefficient, along with
all original prefix/\(\mathfrak k,\mathfrak s,\mathfrak w\) collisions,
fixed numerator zeros and prefix phases.
The intrinsic arithmetic \(H_{\mathfrak p}\) still depends only on
its own factor assignments. Multiplication by these outer factors
preserves its divisor bound. A forbidden prefix is zero before
a row bound is applied.

The remaining long-variable deletion is
\[
 Q=\operatorname{rad}_{\rm good}
    (\mathfrak p\mathfrak n\mathfrak k\mathfrak s\mathfrak w).
\]
It is not a substitute for the first factor in
\eqref{rg:prefix-mask}, and the growing \(\mathfrak n\) is not
part of a fixed character conductor.
On the nonzero support \(\mathfrak p,\mathfrak k\) are squarefree
and coprime, by the original \(\eta(\mathfrak k)^2\) zero.
The physical cubic frequency is \(\mathfrak t=\mathfrak p\mathfrak k\).

The exact deletion formula is
\begin{equation}\label{rg:deletion-identity}
 \begin{split}
 \mathcal F_\psi(T;\mathfrak t;Q,I)
 =\sum_{\substack{\mathfrak h\mid Q\\(\mathfrak h,\mathfrak t)=1}}
 &\frac{\mu(\mathfrak h)\psi(\mathfrak h)
       \overline{G_{\mathfrak h}^{\rm act}}
       \chi_{3,\mathfrak h}(\alpha_{\mathfrak t})}
                 {\sqrt{N\mathfrak h}}\\
 &\cdot\mathcal F_\psi(T/N\mathfrak h;
                              \mathfrak h\mathfrak t;1,I/N\mathfrak h).
 \end{split}
\end{equation}
Insert \(\sum_{\mathfrak h\mid(Q,\mathfrak v)}\mu(\mathfrak h)\)
and set \(\mathfrak v=\mathfrak h\mathfrak a\).
Gauss CRT and reciprocity combine the actual crosses into
\(\chi_{3,\mathfrak a}(\alpha_{\mathfrak h\mathfrak t})\).
This raw entry is zero at \((\mathfrak a,\mathfrak h)>1\)
and encodes the squarefree-product condition. The outer raw entry
kills \((\mathfrak h,\mathfrak t)>1\).
No zero is inverted, including if a numerator exponent becomes
divisible by three. Original interval endpoints divide literally
by \(N\mathfrak h\).

Here \(T=D/N\mathfrak p\), and the exact product is
\begin{equation}\label{rg:deletion-product}
 (T/N\mathfrak h)N(\mathfrak h\mathfrak t)=D N\mathfrak k .
\end{equation}
There is no common-long-divisor decomposition between the two
columns in this proof. The other original column remains fixed.

Using the sharp bound and the original prefix factor gives
\eqref{rg:small-prefix-weight}.
For \(T/N\mathfrak h\in[1/2,1)\) use the actual unit contribution;
smaller arguments contribute zero, with the profile endpoint rule.
Large frequencies are allowed by the proven all-level sharp bound,
with its trivial fallback. Thus
\[
 \sum_{\mathfrak h\mid Q}(N\mathfrak h)^{-3/8}
 \le\tau(Q)\ll_\epsilon(NQ)^\epsilon .
\]
On the exact support,
\(N\mathfrak n\le2D\), \(N\mathfrak s\le2V\), \(N\mathfrak w<64\),
so \(NQ\ll DZV N\mathfrak k\).
Summing \((N\mathfrak p)^{-3/8}\) over \(N\mathfrak p\le Z\)
by ideal counting proves the small-prefix column bound
\begin{equation}\label{rg:small-prefix}
 V_1(W)(DXA_1)^\epsilon
                  D^{1/4}(N\mathfrak k)^{3/8+\epsilon}Z^{5/8}.
\end{equation}
Every prefix zero was retained until that positive estimate.
No coherence or orientation mean has been removed.

\subsection{The exact discrete prefix and the final normalization}
\label{rg:assembly-proof}
Combine the balanced, large-prefix and small-prefix bounds, and put
\begin{equation}\label{rg:balancing}
 Z=D^{2/9},\qquad
 \sqrt{D/Z}=D^{1/4}Z^{5/8}=D^{7/18}.
\end{equation}
The untouched original \(A_U\) column has the elementary bound
\(B_WD^{1/2+\epsilon}\).
Apply it after the left estimate and use the original \(A_1/D\) once.
All finite source/ray/primitive \(\mathfrak e/\mathfrak w\)
coefficient masses stay fixed, and the actual \(\mathfrak s\)
family is counted by \(O(V)\) only at the end.
The scalar \(g_0\) is bounded and is never divided out.

The original finite \(\mathfrak k\) sum is unchanged in the identity.
Its small-prefix positive majorant is
\begin{equation}\label{rg:k-majorant}
 \sum_{\substack{\mathfrak k\ {\rm good,SF}\\
                      N\mathfrak k\le K_{\rm cut}}}
       (N\mathfrak k)^{-13/8+\epsilon}
 \le\zeta_K(13/8-\epsilon),\qquad \epsilon<5/8 .
\end{equation}
Product and divisor losses can be allocated smaller exponents first;
their \(N\mathfrak k\) factors are absorbed by this convergent series.
Balanced and large-prefix pieces use the easier exponent \(2\).
This is a bound on a finite sum, not an Euler replacement of the model.

For the optional original unequal-norm projection, a fixed rational
norm has \(O_\epsilon(D^\epsilon)\) ideal coordinates.
The actual \(A_U\) and every fixed HB component are divisor-bounded.
Removing that same fixed-\(\mathfrak n\) fibre costs
\(O(B_WD^{-1/2+\epsilon})\) in the left column, and therefore
\(O(B_W^2A_1VD^{-1+\epsilon})\) in the normalized form.
Conjugate ideals remain different, and a norm inequality has not
been made multiplicative.

This proves the leading-model part of \eqref{rg:prefix-bound}.
For the original exact prefix, all its Bessel arguments are at most
\(\Lambda=A_1X/D^2\), and
\[
 |\dot G(x)-g_0|\le3C_0^{-3/2}x .
\]
There are \(O(D^2)\) annular column pairs and
\(A_1/(D\sqrt{F_\Omega})\le A_1/D^2\).
The coefficient product is \(O(B_W^2D^\epsilon)\),
the original finite \(\mathfrak k\) absolute budget is bounded by
\(\sum(N\mathfrak k)^{-2}\), and the actual sixth-part band has
\(O(V)\) ideals. Only in this remainder are these absolute counts
taken. They give \eqref{rg:model-error} directly for \(A_U,A_U\).
No \(A_U\)-to-\(A_Y\) error, principal density or scalar ray-count
split is introduced.

Finally the exact powers are
\begin{equation}\label{rg:all-powers}
 \begin{aligned}
 3/4+1/16-1/9&=101/144,&3/4-101/144&=7/144,\\
 3/4+1/16-1/6&=31/48,&
 13/16+(3/4+3/8-2)&=-1/16 .
 \end{aligned}
\end{equation}
The same norm-fibre correction has power \(-3/16\).
This completes the pure finite-prefix theorem with its once-only
normalization and all original finite cutoffs.

\subsection{Verification scope and the remaining arithmetic comparison}
\label{rg:evidence}
The accepted evidence includes exact rational cusp matrices and
finite quotient/coset calculations, angular-coordinate and Mellin
checks, literal primary-support cases, affine coefficient transport,
and synthetic scale-composition/exponent checks.
The typed exponent grid is a formal subfamily; it is not exhaustive
physical dyad coverage after the constants in
\eqref{rg:box-product}. The proof uses the separate exact
\(D\ge2^{64}\) inequalities for those constants.
None of these checks is actual-family energy data, training,
asymptotics or a kernel-checked Lean proof.

All original sources, primary versions, raw outputs and failed
preliminary harnesses remain preserved. The new theorem is a
prose proof with the identified external theta and cubic-norm
inputs. Its domain is the representation-fixed pure
\(\mathcal T_X\) prefix in \eqref{rg:exact-prefix}.
The infinite dual complement, other pure sectors, mixed and
neighboring blocks, short/principal assembly, independent positive
arithmetic determinant and zero-location comparisons remain open.
The Dirichlet, Hecke, Dedekind, elliptic and general Artin scopes
retain their separate hypotheses; no Artin holomorphy, desired
positivity, zero location or GRH premise has been inserted.
The whole signed and global unsigned comparisons have not changed.
